% \iffalse meta-comment
% file: cescg.dtx
% version: 0.1.2
% date: November 25, 2005
%
% Provides `cescg' document class for use with LaTeX2e.
% Formats a document according to CESCG conference proceedings specifications.
%
% \iffalse
%<bib>@String{COPYRIGHT = {
%<*!eps>
% \fi
%% Copyright (c) 1999, 2005
%%               Institute of Computer Graphics and Algorithms
%%               TU Vienna, Austria
%% Based on the LaTeX2e document class SCCG 
%% Copyright (C) 1999 Pavel Chalmoviansky
%%                    Katedra geometrie
%%                    Faculty of Mathematics and Physics
%%                    Comenius University, Bratislava
%%                    chalmo@fmph.uniba.sk
%%
%% This file is distributed in the hope that it will be useful,
%% but WITHOUT ANY WARRANTY; without even the implied warranty of
%% MERCHANTABILITY or FITNESS FOR A PARTICULAR PURPOSE.
%%
% \iffalse
%<bib>}}
% \fi
%
% This file may be freely copied in whole or in part under the following 
%  conditions:
% 1) Unmodified copies may be freely distributed separately or as part
%    of larger packages.
% 2) You may make modified copies of this file for your own personal use,
%    but you may not distribute modified copies. You may place any 
%    modifications in a cescg.cfg file (see the documentation), and distribute 
%    that file individually or along with this source. However, if you do
%    so, the .cfg file must contain code to issue a warning when it is
%    invoked by LaTeX that states:
%       a) The fact that it is being loaded and it is not part of the
%           standard distribution
%       b) What it does (briefly)
%       c) That it can be disabled by removing or renaming the .cfg file
% 3) Excerpts from this file may be freely used and distributed as part
%    of other macro packages.
%
% Portions of this file (namely the \@maketitle macro) are modified versions 
%  of those provided with the standard LaTeX files `article.cls' file and 
%  `size10.clo', and their distribution may be further restricted by the 
%  LaTeX copyright. (Probably not, but better safe than sorry.)
%  
%
% \fi
%
% \CheckSum{848}
% \iffalse
%<bib>@String{CHARACTERTABLE = {
% \fi
%% \CharacterTable
%%  {Upper-case    \A\B\C\D\E\F\G\H\I\J\K\L\M\N\O\P\Q\R\S\T\U\V\W\X\Y\Z
%%   Lower-case    \a\b\c\d\e\f\g\h\i\j\k\l\m\n\o\p\q\r\s\t\u\v\w\x\y\z
%%   Digits        \0\1\2\3\4\5\6\7\8\9
%%   Exclamation   \!     Double quote  \"     Hash (number) \#
%%   Dollar        \$     Percent       \%     Ampersand     \&
%%   Acute accent  \'     Left paren    \(     Right paren   \)
%%   Asterisk      \*     Plus          \+     Comma         \,
%%   Minus         \-     Point         \.     Solidus       \/
%%   Colon         \:     Semicolon     \;     Less than     \<
%%   Equals        \=     Greater than  \>     Question mark \?
%%   Commercial at \@     Left bracket  \[     Backslash     \\
%%   Right bracket \]     Circumflex    \^     Underscore    \_
%%   Grave accent  \`     Left brace    \{     Vertical bar  \|
%%   Right brace   \}     Tilde         \~}
%%
% \iffalse
%</!eps>
%<bib>}}
% \fi
% 
% \iffalse
%    \begin{macrocode}
%<class>\NeedsTeXFormat{LaTeX2e}[1994/06/01]
%<class>\ProvidesClass{cescg}
%<sample>\ProvidesFile{cescgsmpl.tex}
%<*driver>
        \ProvidesFile{cescg.drv}
%</driver>
%<*class|sample|driver>
            [2005/11/29 v0.1.4 
%</class|sample|driver>
%<class> CESCG proceedings document class]
%<sample> CESCG proceedings sample file]
%<*driver> 
       CESCG proceedings document class]
%</driver> 
%    \end{macrocode}
%
%
% \section{Driver for this document}
%
% The following contains the driver for generating this document.
% It can be extracted from the |.dtx| file with the \dst{} program.
%
%    \begin{macrocode}
%<*driver>
\documentclass{ltxdoc}
\GetFileInfo{cescg.drv}
%\DisableCrossrefs
\RecordChanges
\OnlyDescription % comment out for full source description
\begin{document}
    \DocInput{cescg.dtx}
\end{document}
%</driver>
%    \end{macrocode}
%\fi
%
% \newenvironment{mylist}{%
%       \begin{list}{\labelitemi}{%
%               \addtolength{\leftmargin}{2em}%
%               \setlength{\itemindent}{-1em}%
%               \setlength{\parsep}{0in}%
%               \setlength{\itemsep}{0.5ex}}}{\end{list}}
%
% \changes{v0.1.0}{1999/1/31}{TT: Created the CESCG document class,
%                                based on Pavel Chalmoviansky's \texttt{sccg.cls},
%                                which is itself based on the version 1.2.1 of `siggraph.dtx'.}
%
% \changes{v0.1.1}{2005/11/12}{MW:
%                  (1) cls-file: pass option a4paper to article class automatically
%                  (2) template file:
%                      (a) changed to \texttt{graphicx} graphic driver instead of \texttt{epsfig}
%                      (b) included instructions to include \texttt{.png} files directly
%                      (c) included instructions to include \texttt{.eps} files for \texttt{pdflatex}}
% 
% \changes{v0.1.2}{2005/11/25}{GZ: fixed \texttt{.dtx} file, updated documentation}
% \changes{v0.1.3}{2005/11/28}{GZ: Bugfix, removed buggy option epstopdf}
% \changes{v0.1.4}{2010/02/18}{MI: Footer line}
%
% \title{The |CESCG| Document Class Users' Guide\thanks{
%    This file has version number \fileversion{} dated \filedate{}.
%    The most recent version can be found at
%    \texttt{http://www.cg.tuwien.ac.at/studentwork/CESCG/}}}
%
% \author{Thomas Theu{\ss}l, Michael Wimmer, Georg Zotti, Martin Ilcik\thanks{\{theussl,wimmer,gzotti,ilcik\}@cg.tuwien.ac.at}}
%
% \date{\filedate}
%
% \maketitle
%
% \begin{abstract}
%   This document class modifies the standard |article| class to conform to
%   the specifications for papers published in the proceedings of the annual
%   Spring Conference on Computer Graphics.
%   It sets all the necessary layout parameters and handles the differences
%   in format between a paper being submitted for blind review and a camera
%   ready copy of an accepted paper. (currently discouraged.)
%   Several additional features are also provided.
%
%   \vspace{\baselineskip}
%   \begin{quote}\scshape
%     \textbf{DISCLAIMER}:  This class is distributed in the hope that it will be useful,
%     and to promote compatibility of documents created at different sites.
%     However it is provided \textbf{WITHOUT ANY WARRANTY}; without even the implied
%     warranty of \textbf{MERCHANTABILITY} or \textbf{FITNESS FOR A PARTICULAR PURPOSE}.
%     Please see the |cescg.dtx| source file for copyright and distribution
%     information.
%   \end{quote}
% \end{abstract}
%
% \section{Introduction}
%
% This document describes the |cescg| document class, for use with \LaTeXe{}.
% It loads in the |article| class and modifies several parameters
% to comply with the layout specifications for a paper submitted to
% the CESCG conference.
% Currently, these specifications are:
% \begin{mylist}
%   \label{layoutlist}
%   \item A4 paper size
%   \item $2$ cm margins on top and sides, $2.6$ margin on bottom
%   \item two column mode with $0.7$ cm column separation
%   \item 10 point type on 12 point baselines
%   \item no page numbers
% \end{mylist}
% Also set by the class but not part of the CESCG specifications are:
% \begin{mylist}
%   \item titles and section headings are bold sans serif
%   \item 1 em paragraph indentation
%   \item |flushbottom| mode
%   \item whole title flushed left or right
% \end{mylist}
%
% The class also allows the user to specify whether the paper is being
% submitted for blind review or as a final, camera-ready document.
% When review mode is specified, it will automatically generate a
% cover sheet for the paper, modify the title block on the first page
% to protect anonymity, and (optionally) switch to a more widely spaced
% layout.
% Mechanisms are also provided to allow the user to introduce text conditional
% on the chosen submission mode.
%
% \section{Class Usage}
%
% To use this class, specify
%  \begin{quote}
%     |\documentclass[|\meta{options}|]{cescg}|
%  \end{quote}
%  at the beginning of your document.
% The following options may be included in the |\documentclass| command:
%
% \begin{mylist}
%  \item \DescribeMacro{cameraready}|cameraready|,
%        \DescribeMacro{preprint}|preprint|,
%        \DescribeMacro{review}|review|, or
%        \DescribeMacro{widereview}|widereview|
%       \newline These specify the submission mode for the paper.
%       |cameraready|, the default, sets the layout parameters as
%       described in section~\ref{layoutlist}.
%       The other options are deprecated and may go away, please don't use them.
% \iffalse       % GZ, 2005/11/25
%       The |preprint| option adds page numbers and prints the message
%       ``To appear in the CESCG conference proceedings'' in the top
%       margin.  This message can be changed with the
%       \DescribeMacro{\preprinttext}|\preprinttext{|\meta{message}|}|
%       command.
%       The |review| option uses the same layout as |cameraready| but
%       adds page numbers and modifies the title block as described
%       in the following sections, and causes a cover page to
%       be generated.
%       The |widereview| option is for people who prefer review papers
%       that have a more widely spaced layout.
%       It uses one-column mode with 12 point, double-spaced type.
% \fi
%  \item \DescribeMacro{flushtitleleft}|flushtitleleft| or
%        \DescribeMacro{flushtitleright}|flushtitleright|
%       \newline These option specify, where the title will be flushed.
%       The default option is |flushright|.
%  \item \DescribeMacro{onecolumn}|onecolumn| or
%        \DescribeMacro{twocolumn}|twocolumn|
%  \item \DescribeMacro{singlespace}|singlespace| or
%        \DescribeMacro{doublespace}|doublespace|
%  \item \DescribeMacro{10pt}|10pt|,
%        \DescribeMacro{11pt}|11pt|, or
%        \DescribeMacro{12pt}|12pt|
%       \newline These three sets of options will override the default
%       settings for the number of columns, line spacing, and point size
%       determined by the submission mode.
%       Overriding the settings when in camera-ready mode, while allowed,
%       is a violation of the submission specifications and will produce
%       a warning message.
%   \item \DescribeMacro{epstopdf}|epstopdf| loads the package of the same name,
%         allowing the inclusion of |.eps| figures when using pdf\LaTeX.
%
%  \item Additionally, any other options supported by the |article| class
%       may be used, except for |landscape| and |titlepage|, which are
%       disabled.
% \end{mylist}
%       
%
% \section{Title Block, Abstract, and Cover Page}
%
% \subsection{Data fields}
%
% To produce the title block and cover sheet, a number of data fields must
%  be set in the preamble of the paper (i.e.~before |\begin{document}|).
% Each of the following commands takes a single argument, the value to which
%  the corresponding field should be set.
% Footnotes to any of these items should be specified with the |\thanks|
%  command, rather than |\footnote|.
%
% \begin{mylist}
%   \item \DescribeMacro{\title}|\title|
%
%       The title of the paper.
%
%   \item \DescribeMacro{\author}|\author|
%
%       The author(s) of the paper and their individual affiliations.
%       (See the sample file for examples.)
%
%   \item \DescribeMacro{\affiliation}|\affiliation|
%
%       If all the authors share an affiliation, it can be set with this
%       command rather than individually within the author command.
%       The affiliation will appear centered below the author list.
%   \item \DescribeMacro{\supervisor}|\supervisor|
%
%		This command is used to set the name of the supervisor.
%		It will appear centered below the authors listing.
%
% \iffalse GZ: 2005-11-28: taken out, we don't need it currently.
%<*gobble>
%   \item \DescribeMacro{\category}|\category| and
%         \DescribeMacro{\format}|\format|
%
%       The paper category (research, systems, application, survey, \ldots) and format (print, video).
%
%   \item \DescribeMacro{\contactname}|\contactname|,
%         \DescribeMacro{\contactaddress}|\contactaddress|,
%         \DescribeMacro{\contactphone}|\contactphone|,
%         \DescribeMacro{\contactfax}|\contactfax|,
%         \DescribeMacro{\contactemail}|\contactemail|
%
%       These should be used to set the information about the person to
%       be contacted about the paper. They will be shown only on cover page.
%</gobble>
% \fi
%   \item \DescribeMacro{\keywords}|\keywords|
%
%       A list of keywords to be printed on the cover page
%
% \iffalse GZ: 2005-11-28: taken out, we don't need it currently.
%<*gobble>
%   \item \DescribeMacro{\estpages}|\estpages|
%
%       The estimated number of pages that the paper would require in the
%       conference proceedings.
%</gobble>
% \fi
% \end{mylist}
%
% \subsection{Title block}
%
% The title block should be generated with the normal \LaTeX{}
%  \DescribeMacro{\maketitle}|\maketitle| command, place after
%  the |\begin{document}| command.
% In camera ready or preprint mode it will contain:
% \begin{mylist}
%   \item Title
%   \item Authors
%   \item Affiliation (if any)
% \end{mylist}
% For a review paper it will instead contain:
% \begin{mylist}
%   \item Title
%   \item Paper category
% \end{mylist}
% The review options will also suppress printing of any information
% specified with the |\thanks| command.
%
% \subsection{Abstract}
%
% The abstract should be specified with the |abstract| environment.
% In addition to appearing in the paper itself, its contents will stored
% and will appear on the cover page (in review mode).
%
% \subsection{Keywords}
%
% Authors may wish to include a list of keywords after the abstract.
% The \DescribeMacro{\keywordlist}|\keywordlist| command will include the
%  list set by the |\keywords| command described above.
%
% \iffalse GZ: 2005-11-28: taken out, we don't need it currently.
%<*gobble>
% \subsection{Cover page}
%
% When the |review| or |widereview| option is specified,
% a cover sheet will automatically be generated at the end of the document.
% Placing the \DescribeMacro{\suppresscover}|\suppresscover| command
% in the preamble will prevent this.
% The cover sheet will contain:
% \begin{mylist}
%   \item Title
%   \item Authors
%   \item Affiliation (if any)
%   \item Paper category
%   \item Contact information
%   \item A copy of the abstract
%   \item List of keywords
%   \item Estimated number of pages
% \end{mylist}
%</gobble>
% \fi
%
% \section{Other Features}
%
% \subsection{Title and section heading font}
%
% By default, this class uses bold sans-serif for the title and
% section heading font.
% This behavior can be altered with the
% \DescribeMacro{\sectionfont}|\sectionfont{|\meta{fontcommands}|}| command.
% \meta{fontcommands} should be one or more font changing commands,
% for instance ``|\rmfamily\scshape|'' or ``|\fontfamily{pcr}\selectfont|''.
% Normally, you should not use this command, though, as it breaks the
% formatting guidelines.
% Please do \emph{not} load packages that change the regular text font,               
% e.g. \texttt{times} etc., for the same reason.  
%
% \subsection{Conditional text}
%
% \DescribeMacro{\ifcamera}\DescribeMacro{\ifreview}
% \DescribeMacro{\ifcameraelse}\DescribeMacro{\ifreviewelse}
% The class also provides the following four commands for writing text
% conditional on whether camera-ready or review mode has been
% chosen:
% \begin{quote}
%       |\ifcamera{|\meta{thenclause}|}|
%
%       |\ifreview{|\meta{thenclause}|}|
%
%       |\ifcameraelse{|\meta{thenclause}|}{|\meta{elseclause}|}|
%
%       |\ifreviewelse{|\meta{thenclause}|}{|\meta{elseclause}|}|
% \end{quote}
% The first two will generate the \meta{thenclause} if camera-ready
% (resp. review) mode is chosen, otherwise they resolve to a null string.
% The second two are the same as the first two but generate the
% \meta{elseclause} instead of a null string if the condition is not met.
%
% These can be used for writing the paper so as to protect anonymity for
% a submission for review, without having to worry about going back over it
% months later to put the identifying information back in if the paper
% is accepted.  For example:
% \begin{quote}^^A
%   |In \cite{FooBar1994}, \ifcameraelse{we}{Foo and Bar} showed that ...|
% \end{quote}
% or
% \begin{quote}^^A
%   |as illustrated in the following \ifcameraelse{three}{two} images:|
%
%   [include generic image]
%
%   |\ifcamera{| [include our logo] |}|
%
%   [include another generic image]
% \end{quote}
%
% \subsection{Configuration file}
%
% Individual users or sites may wish to modify the behavior of this
% class (for instance, to use a particular set of fonts.)
% Rather than having people hack the |.cls| file, which can lead to
% various incompatible versions at different sites, the class provides
% for a configuration file.
% After the class has finished loading, it will check to see if a
% |cescg.cfg| file exists in the \TeX{} search path.
% If so, it will be input.
% Any user-- or site--specific modifications should be placed in this file.
% \textit{(See the copyright notice in the |cescg.dtx| source file for
% information on distribution of configuration files.)}
%
% \section{Example File}
%
%  This class comes with a file named ``cescgsmpl.tex'', which provides
%   a simple example of how to use this class.
%  To compile it, run
%  \begin{verbatim}
%    latex cescgsmpl
%    bibtex cescgsmpl
%    latex cescgsmpl
%    latex cescgsmpl
%  \end{verbatim}
%
% \StopEventually{\PrintChanges}
%
% \iffalse
%<*class>
% \fi
%
% \section{Implementation}
%
% \subsection{Initial stuff}
%
%  \begin{macro}{\ifcescg@camera}
%  \begin{macro}{\ifcescg@preprint}
%  \begin{macro}{\ifcescg@review}
%  \begin{macro}{\cescg@columnmode}
%  \begin{macro}{\cescg@pointsize}
%  \begin{macro}{\ifcescg@doublespaced}
%   First off, we need to initialize a few variables.
%   |cescg@camera| is a boolean which determines whether or not we're in
%   camera ready mode.
%   |\cescg@preprint|, |\cescg@columnmode|, |\cescg@pointsize|,
%   |cescg@doublespaced| and |flushtitleleft| are pretty obvious.
%    \begin{macrocode}
\newif\ifcescg@camera
\newif\ifcescg@preprint
\newif\ifcescg@review
\newcommand{\cescg@columnmode}{}
\newcommand{\cescg@pointsize}{}
\newif\ifcescg@doublespaced
\newif\ifcescg@flushtitleleft
%    \end{macrocode}
%  \end{macro}
%  \end{macro}
%  \end{macro}
%  \end{macro}
%  \end{macro}
%  \end{macro}
%
% \subsection{Option processing}
%
% \subsubsection{Submission mode}
%
%  \begin{macro}{cameraready}
%  \begin{macro}{preprint}
%  \begin{macro}{review}
%  \begin{macro}{widereview}
%   The |cameraready|, |preprint|, |review|, and |widereview| options set the 
%   variables defined above appropriately.
%    \begin{macrocode}
\DeclareOption{cameraready}{%
    \cescg@cameratrue%
    \cescg@preprintfalse%
    \cescg@reviewfalse%
    \renewcommand{\cescg@columnmode}{twocolumn}%
    \cescg@doublespacedfalse%
    \renewcommand{\cescg@pointsize}{10pt}}
\DeclareOption{review}{%
    \cescg@camerafalse%
    \cescg@preprintfalse%
    \cescg@reviewtrue%
    \renewcommand{\cescg@columnmode}{twocolumn}%
    \cescg@doublespacedfalse%
    \renewcommand{\cescg@pointsize}{10pt}}
\DeclareOption{widereview}{%
    \cescg@camerafalse%
    \cescg@preprintfalse%
    \cescg@reviewtrue%
    \renewcommand{\cescg@columnmode}{onecolumn}%
    \cescg@doublespacedtrue%
    \renewcommand{\cescg@pointsize}{12pt}}
\DeclareOption{preprint}{%
    \cescg@camerafalse%
    \cescg@preprinttrue%
    \cescg@reviewfalse%
    \renewcommand{\cescg@columnmode}{twocolumn}%
    \cescg@doublespacedfalse%
    \renewcommand{\cescg@pointsize}{10pt}}
%    \end{macrocode}
%  \end{macro}
%  \end{macro}
%  \end{macro}
%  \end{macro}
%  By default, |cameraready| mode should be used.
%    \begin{macrocode}
\ExecuteOptions{cameraready}
%    \end{macrocode}
%
% \subsubsection{Layout parameters}
%
%  \begin{macro}{\cescg@optwarning}
%   The |\cescg@optwarning| command produces a warning message when one
%   tries to override the default layout parameters while in camera-ready
%   mode.
%    \begin{macrocode}
\newcommand{\cescg@optwarning}[1]{%
    \ifcescg@camera
        \ClassWarningNoLine{cescg}%
            {#1 option used in camera-ready mode.\MessageBreak
             This violates submission specifications}
    \fi
}
%    \end{macrocode}
%  \end{macro}
%
%  \begin{macro}{onecolumn}
%  \begin{macro}{twocolumn}
%  \begin{macro}{10pt}
%  \begin{macro}{11pt}
%  \begin{macro}{12pt}
%  \begin{macro}{doublespace}
%  \begin{macro}{singlespace}
%  \begin{macro}{flushtitleft}
%  \begin{macro}{flushtitleright}
%   The various layout parameters can be set individually, overriding the
%   values set by the specified submission mode.
%    \begin{macrocode}
\DeclareOption{onecolumn}{%
    \cescg@optwarning{onecolumn}
    \renewcommand{\cescg@columnmode}{onecolumn}}
\DeclareOption{twocolumn}{%
    \renewcommand{\cescg@columnmode}{twocolumn}}
\DeclareOption{10pt}{%
    \cescg@optwarning{10pt}
    \renewcommand{\cescg@pointsize}{10pt}}
\DeclareOption{11pt}{%
    \cescg@optwarning{11pt}
    \renewcommand{\cescg@pointsize}{11pt}}
\DeclareOption{12pt}{%
    \cescg@optwarning{12pt}
    \renewcommand{\cescg@pointsize}{12pt}}
\DeclareOption{singlespace}{%
    \cescg@doublespacedfalse}
\DeclareOption{doublespace}{%
    \cescg@optwarning{doublespace}
    \cescg@doublespacedtrue}
\DeclareOption{flushtitleleft}{%
    \cescg@flushtitlelefttrue}
\DeclareOption{flushtitleright}{%
    \cescg@flushtitleleftfalse}
%    \end{macrocode}
%  \end{macro}
%  \end{macro}
%  \end{macro}
%  \end{macro}
%  \end{macro}
%  \end{macro}
%  \end{macro}
%  \end{macro}
%   Set the default flushtitle to right.
%    \begin{macrocode}
\ExecuteOptions{flushtitleright}
%    \end{macrocode}
%
% \subsubsection{Remaining options}
%
%  Disable the |titlepage| and |landscape| options, and pass any remaining
%  options on to the |article| class.
%
%    \begin{macrocode}
\DeclareOption{titlepage}{%
    \OptionNotUsed%
    \ClassWarningNoLine{cescg}{titlepage option not allowed}}
\DeclareOption{landscape}{%
    \OptionNotUsed%
    \ClassWarningNoLine{cescg}{landscape option not allowed}}
\DeclareOption*{\PassOptionsToClass{\CurrentOption}{article}}
%    \end{macrocode}
%
% \subsubsection{Option processing}
%
%  Process the options, cause the appropriate column and point size options
%  to be passed to the |article| class, and load the |article| class.
%
%    \begin{macrocode}
\ProcessOptions
\PassOptionsToClass{a4paper}{article}
\PassOptionsToClass{\cescg@columnmode}{article}
\ifdim\cescg@pointsize>10pt
    \PassOptionsToClass{\cescg@pointsize}{article}
\fi
\LoadClass{article}
% \changes{v0.1.2}{2005/11/25}{GZ: load font packages} 
\RequirePackage{mathptmx} 
\RequirePackage[scaled=.90]{helvet} 
\RequirePackage{courier} 
%    \end{macrocode}
%
% \subsection{Preprint text}
%  \begin{macro}{\preprinttext}
%   Set up default message for preprints and command to change it.
%    \begin{macrocode}
\newcommand{\cescg@preprinttext}{To appear in the CESCG conference proceedings}
\newcommand{\preprinttext}[1]{\renewcommand{\cescg@preprinttext}{#1}}
\newcommand{\cescg@footertext}{Proceedings of CESCG 2013: The 17th Central European Seminar on Computer Graphics (non-peer-reviewed)}
\newcommand{\footertext}[1]{\renewcommand{\cescg@footertext}{#1}}
%    \end{macrocode}
%  \end{macro}
%
%
% \subsection{Conditional text}
%
%  \begin{macro}{\ifcamera}
%  \begin{macro}{\ifreview}
%  \begin{macro}{\ifcameraelse}
%  \begin{macro}{\ifreviewelse}
%   Here we define the commands to include text conditionally, based on
%   the submission mode.
%    \begin{macrocode}
\newcommand{\ifcamera}[1]{\ifcescg@review \else #1 \fi}
\newcommand{\ifreview}[1]{\ifcescg@review #1 \fi}
\newcommand{\ifcameraelse}[2]{\ifcescg@review #2 \else #1 \fi}
\newcommand{\ifreviewelse}[2]{\ifcescg@review #1 \else #2 \fi}
%    \end{macrocode}
%  \end{macro}
%  \end{macro}
%  \end{macro}
%  \end{macro}
%
% \subsection{Page layout}
%
% Set the margins, header and footer sizes, column separation, and indentation.
%
%    \begin{macrocode}
\setlength{\textheight}{24cm}
\setlength{\topmargin}{-0.85cm}
\setlength{\headheight}{0.5cm}
\setlength{\headsep}{0.65cm}
\setlength{\footskip}{1.25cm}
\flushbottom
\setlength{\hoffset}{-0.54cm}
\setlength{\textwidth}{17cm}
\setlength{\oddsidemargin}{0cm}
\setlength{\evensidemargin}{0cm}
\setlength{\marginparsep}{0cm}
\setlength{\marginparwidth}{0pt}
\setlength{\columnsep}{0.7cm}
\setlength{\parindent}{1em}
%    \end{macrocode}
%
% Modify plain and empty page styles to turn off page numbering or turn on
%  preprint message.
%
%    \begin{macrocode}
\renewcommand{\ps@plain}%
    {%
    \ifcescg@preprint%
\renewcommand{\@oddhead}{\hfil\textit{\large\cescg@preprinttext}\hfil}%
    \else%
\renewcommand{\@oddhead}{}%
    \fi%
    \ifcescg@camera%
    \renewcommand{\@oddfoot}{\hfil\textit{\cescg@footertext}\hfil}%
    \else%
\renewcommand{\@oddfoot}{\hfil\textrm{\thepage}\hfil}%
    \fi%	
    \renewcommand{\@evenhead}{\@oddhead}%
    \renewcommand{\@evenfoot}{\@oddfoot}%
    }
\renewcommand{\ps@empty}%
    {%
    \ifcescg@preprint%
\renewcommand{\@oddhead}{\hfil\textit{\large\cescg@preprinttext}\hfil}%
    \else%
    \renewcommand{\@oddhead}{}%
    \fi%
    \renewcommand{\@oddfoot}{\hfil\textit{\cescg@footertext}\hfil}%
    \renewcommand{\@evenhead}{\@oddhead}%
    \renewcommand{\@evenfoot}{\@oddfoot}%
    }
\pagestyle{plain}
%    \end{macrocode}
%
% \subsection{Section headings}
%
%  \begin{macro}{\sectionfont}
%   The |\sectionfont| command lets the user modify the font used for
%   section headings and the title.
%    \begin{macrocode}
\newcommand{\cescg@sectionfont}{}
\newcommand{\sectionfont}[1]{\renewcommand{\cescg@sectionfont}{#1}}
%    \end{macrocode}
%  \end{macro}
%
%  \begin{macro}{\section}
%  \begin{macro}{\subsection}
%  \begin{macro}{\subsubsection}
%  \begin{macro}{\paragraph}
%  \begin{macro}{\subparagraph}
%   We redefine the section headings to use bold sans serif type plus
%   any modifications set with |\sectionfont|.  (Otherwise
%   these are the same definitions given in the |classes.dtx| file.)
%    \begin{macrocode}
\renewcommand\section{\@startsection {section}{1}{\z@}%
                {-3.5ex \@plus -1ex \@minus -.2ex}%
                {2.3ex \@plus.2ex}%
                {\reset@font\Large\sffamily\bfseries\cescg@sectionfont}}
\renewcommand\subsection{\@startsection{subsection}{2}{\z@}%
                {-3.25ex\@plus -1ex \@minus -.2ex}%
                {1.5ex \@plus .2ex}%
                {\reset@font\large\sffamily\bfseries\cescg@sectionfont}}
\renewcommand\subsubsection{\@startsection{subsubsection}{3}{\z@}%
                {-3.25ex\@plus -1ex \@minus -.2ex}%
                {1.5ex \@plus .2ex}%
                {\reset@font\normalsize\sffamily\bfseries\cescg@sectionfont}}
\renewcommand\paragraph{\@startsection{paragraph}{4}{\z@}%
                {3.25ex \@plus1ex \@minus.2ex}%
                {-1em}%
                {\reset@font\normalsize\sffamily\bfseries\cescg@sectionfont}}
\renewcommand\subparagraph{\@startsection{subparagraph}{5}{\parindent}%
                {3.25ex \@plus1ex \@minus .2ex}%
                {-1em}%
                {\reset@font\normalsize\sffamily\bfseries\cescg@sectionfont}}
%    \end{macrocode}
%  \end{macro}
%  \end{macro}
%  \end{macro}
%  \end{macro}
%  \end{macro}
%
%
%
% \subsection{Title block}
%
%  \begin{macro}{\cescg@empty}
%   This is used to check whether the affiliation and teaser have been set.
%    \begin{macrocode}
\newcommand{\cescg@empty}{}
%    \end{macrocode}
%  \end{macro}
%
%  \begin{macro}{\affiliation}
%  \begin{macro}{\cescg@affiliation}
%   We define the |\affiliation| macro to set the authors' affiliation,
%   and the |\cescg@affiliation| macro to hold it.
%    \begin{macrocode}
\newcommand{\cescg@affiliation}{}
\newcommand{\affiliation}[1]{\renewcommand{\cescg@affiliation}{#1}}
%    \end{macrocode}
%  \end{macro}
%  \end{macro}
%
%  \begin{macro}{\supervisor}
%  \begin{macro}{\cescg@supervisor}
%   We define the |\supervisor| macro to set the supervisor name,
%   and the |\cescg@supervisor| macro to hold it.
%    \begin{macrocode}
\newcommand{\cescg@supervisor}{}
\newcommand{\supervisor}[1]{\renewcommand{\cescg@supervisor}{\emph{Supervised by: {#1}}}}
%    \end{macrocode}
%  \end{macro}
%  \end{macro}
%
%  \begin{macro}{\category}
%  \begin{macro}{\cescg@category}
%  \begin{macro}{\format}
%  \begin{macro}{\cescg@format}
%   These macros set and hold the paper's category and format.
%    \begin{macrocode}
\newcommand{\cescg@category}{}
\newcommand{\category}[1]{\renewcommand{\cescg@category}{#1}}
\newcommand{\cescgcategory}[1]{\category{#1}}
\newcommand{\cescg@format}{}
\newcommand{\format}[1]{\renewcommand{\cescg@format}{#1}}
\newcommand{\cescgformat}[1]{\format{#1}}
%    \end{macrocode}
%  \end{macro}
%  \end{macro}
%  \end{macro}
%  \end{macro}
%
%  \begin{macro}{\titlespace}
%   We also need to set up length for the distance between the title block
%   and the rest of the text.  We initialize it to $\frac{1}{4}$ inches.
%    \begin{macrocode}
\newlength{\titlespace}
\setlength{\titlespace}{0.25in}
%    \end{macrocode}
%  \end{macro}
%  \end{macro}
%
%  \begin{macro}{\@maketitle}
%   Now we redefine the \@maketitle command to properly format the title block.
%   For camera-ready papers, it should have the title, authors, and affiliation.
%   For review papers, it should have the title, paper category and format,
%   plus all information specified with the |\thanks| command should be
%   suppressed.
%    \begin{macrocode}
\ifcescg@review
    \renewcommand{\@maketitle}{%
        \begin{flushright}%
            \renewcommand{\thanks}[1]{}
            {\LARGE\sffamily\bfseries\cescg@sectionfont \@title \par}%
            \vspace{1\baselineskip}%
            {Category: \cescg@category \par}%
            \vspace{\titlespace}%
        \end{flushright} \par%
    }
\else
    \renewcommand{\@maketitle}{%
        \begin{center}%
            {\LARGE\sffamily\bfseries\cescg@sectionfont \@title \par}%
            \vspace{1\baselineskip}%
            \large
            \begin{tabular}[t]{@{}c@{}}%
              \@author
            \end{tabular}\par%
			\cescg@supervisor\par%
            \ifx\cescg@affiliation\cescg@empty \else%
                \vspace{1\baselineskip}%
                \cescg@affiliation\par%
            \fi%
            \vspace{\titlespace}%
        \end{center}%
        \par%
    }
\fi
%    \end{macrocode}
%  \end{macro}
%
%  \begin{macro}{\maketitle}
%   We also need to make several modifications to the |\maketitle| command.
%   Since |\maketitle| wipes out the values of the title, author, etc\ldots,
%   we have to make sure a copy gets saved.
%   We also save a copy of the definition of |\and|, since |\maketitle|
%   clears it and we will still need it for the cover page
%   After the title block is output,
%   the page numbering for the first page should be turned off, and
%   double spacing, if selected, should be turned on.
%    \begin{macrocode}
\let\cescg@origmaketitle\maketitle%
\let\and\crcr%
\let\cescg@origand\and%
\renewcommand{\maketitle}{%
    \let\cescg@title\@title%
    \let\cescg@author\@author%
    \cescg@origmaketitle%
    \thispagestyle{empty}%
    \ifcescg@doublespaced%
        \renewcommand{\baselinestretch}{1.66}\selectfont%
    \fi
}
%    \end{macrocode}
%  \end{macro}
%
%
%
% \subsection{Capturing the abstract}
%
%  \begin{environment}{abstract}
%   In order to have the abstract appear on the cover sheet, we need to
%   redefine the |abstract| environment to store its contents so that
%   it can be used later.
%   This seemingly trivial problem is actually quite difficult.
%   The code to do it is a modified version of the |\@collectbody| stuff
%   from the AMS-\LaTeX{} source code.
%   See that code for a more detailed explanation of how it works.
%
%  \begin{macro}{\cescg@abs}
%   The |\cescg@abs| token list is used to store the text of the abstract.
%    \begin{macrocode}
\newtoks\cescg@abs
%    \end{macrocode}
%  \end{macro}
%
%   We only make these modifications if we are in review mode.
%    \begin{macrocode}
\ifcescg@review
%    \end{macrocode}
%
%  \begin{macro}{\cescg@add}
%   The |\cescg@add| command appends its argument to |\cescg@abs|.
%    \begin{macrocode}
    \long\def\cescg@add#1{\global\cescg@abs\expandafter{\the\cescg@abs#1}}
%    \end{macrocode}
%  \end{macro}
%
%  \begin{macro}{\cescg@collect}
%  \begin{macro}{\cescg@@collect}
%   The |\cescg@collect| and |\cescg@@collect| commands scan the input for the
%   |\end{abstract}| command and place all input tokens leading up to it
%   into |\cescg@abs|.
%    \begin{macrocode}
    \long\def\cescg@collect{%
        \global\cescg@abs{}%
        \let\abstract\cescg@@collect
        \abstract
    }
    \long\def\cescg@@collect#1\end#2{%
        \def\@tempa{#2}%
        \ifx\@tempa\@currenvir
            \cescg@add{#1}%
            \edef\abstract{\noexpand\end{\@tempa}}%
        \else
            \cescg@add{#1\end{#2}}%
        \fi
        \abstract
    }
%    \end{macrocode}
%  \end{macro}
%  \end{macro}
%
%  \begin{macro}{\cescg@modify}
%   The |\cescg@modify| command makes the actual modifications to the |abstract|
%   environment.
%   The new |abstract| stores its contents with the above macros, then
%   outputs the abstract using its original definition.
%    \begin{macrocode}
    \newcommand{\cescg@modify}{%
        \let\cescg@origabs\abstract%
        \let\cescg@origendabs\endabstract%
        \renewenvironment{abstract}%
            {\cescg@collect}%
            {\begingroup
                \let\abstract\cescg@origabs
                \let\endabstract\cescg@origendabs
                \begin{abstract} \the\cescg@abs \end{abstract}
            \endgroup}%
    }
%    \end{macrocode}
%  \end{macro}
%
%  We want to allow the user to redefine the appearance of the |abstract|
%  environment however he/she likes, so we hold off actually making the 
%  modifications to |abstract| until after the preamble.
%    \begin{macrocode}
    \AtBeginDocument{\cescg@modify}
%    \end{macrocode}
%
%  End of review mode conditional:
%    \begin{macrocode}
\fi
%    \end{macrocode}
%  \end{environment}
%
%
%
% \subsection{Keywords}
%
%  \begin{macro}{\keywordlist}
%   For the keywordlist, we start an unindented paragraph and print
%    ``Keywords:'' in bold, followed by the list set with the |\keywords|
%    macro.
%    \begin{macrocode}
\newcommand{\keywordlist}{%
    \vspace{0.5\baselineskip}%
    \par\noindent \textbf{Keywords: } \cescg@keywords%
}
%    \end{macrocode}
%  \end{macro}
%
% \subsection{Cover page}
% 
%  \begin{macro}{\contactname}
%  \begin{macro}{\cescg@contactname}
%  \begin{macro}{\contactaddress}
%  \begin{macro}{\cescg@contactaddress}
%  \begin{macro}{\contactphone}
%  \begin{macro}{\cescg@contactphone}
%  \begin{macro}{\contactfax}
%  \begin{macro}{\cescg@contactfax}
%  \begin{macro}{\contactemail}
%  \begin{macro}{\cescg@contactemail}
%  \begin{macro}{\keywords}
%  \begin{macro}{\cescg@keywords}
%  \begin{macro}{\estpages}
%  \begin{macro}{\cescg@estpages}
%   We define macros to set and hold all the information about the person
%   to be contacted about the paper and the author to receive compensation.
%    \begin{macrocode}
\newcommand{\cescg@contactname}{}
\newcommand{\contactname}[1]{\renewcommand{\cescg@contactname}{#1}}
\newcommand{\cescg@contactaddress}{}
\newcommand{\contactaddress}[1]{\renewcommand{\cescg@contactaddress}{#1}}
\newcommand{\cescg@contactphone}{}
\newcommand{\contactphone}[1]{\renewcommand{\cescg@contactphone}{#1}}
\newcommand{\cescg@contactfax}{}
\newcommand{\contactfax}[1]{\renewcommand{\cescg@contactfax}{#1}}
\newcommand{\cescg@contactemail}{}
\newcommand{\contactemail}[1]{\renewcommand{\cescg@contactemail}{#1}}
\newcommand{\cescg@keywords}{}
\newcommand{\keywords}[1]{\renewcommand{\cescg@keywords}{#1}}
\newcommand{\cescg@estpages}{}
\newcommand{\estpages}[1]{\renewcommand{\cescg@estpages}{#1}}
%    \end{macrocode}
%  \end{macro}
%  \end{macro}
%  \end{macro}
%  \end{macro}
%  \end{macro}
%  \end{macro}
%  \end{macro}
%  \end{macro}
%  \end{macro}
%  \end{macro}
%  \end{macro}
%  \end{macro}
%  \end{macro}
%  \end{macro}
%  \end{macro}
%
%  \begin{macro}{cescg@cover}
%   Set up a control to determine whether or not to print the cover page.
%    \begin{macrocode}
\newif\ifcescg@cover
\ifcescg@review
    \cescg@covertrue
\else
    \cescg@coverfalse
\fi
\newcommand{\suppresscover}{\cescg@coverfalse}
%    \end{macrocode}
%  \end{macro}
%
%  \begin{macro}{\cescg@coverpage}
%   The |\cescg@coverpage| macro handles the formatting of the cover page.
%    \begin{macrocode}
\newcommand{\cescg@coverpage}{%
    \begin{titlepage}%
        \def\thanks##1{}
        \let\and\cescg@origand
        \normalsize
        \renewcommand{\baselinestretch}{1}\selectfont
        \begin{center}%
            \vspace*{\fill}
            {\LARGE\sffamily\bfseries\cescg@sectionfont \cescg@title \par}%
            \vspace{2\baselineskip}%
            {\large
                \begin{tabular}[t]{c}%
                \cescg@author
             \end{tabular}\par%
            }%			
            \vspace{1\baselineskip}%
			{\large \cescg@supervisor \par}%
			\vspace{1\baselineskip}%
            {\large \cescg@affiliation \par}%
            \addvspace{3\baselineskip}%
            {Category: \cescg@category \par}%
            \vspace{0.5\baselineskip}%
            {Format: \cescg@format \par}%
            \vspace{3\baselineskip}%
            \begin{tabular}{ll}
                Contact: & \cescg@contactname \\[1\baselineskip]
                         & \begin{tabular}[b]{@{}l@{}}
                               \cescg@contactaddress
                               \end{tabular} \\[1\baselineskip]
                phone: & \cescg@contactphone \\
                fax: & \cescg@contactfax \\
                email: & \cescg@contactemail
            \end{tabular}\par%
            \vspace*{3\baselineskip}%
            {Estimated \# of pages: \cescg@estpages \par}%
            \vspace*{\baselineskip}%
            {Keywords: \cescg@keywords \par}%
            \vspace*{\fill}%
            \begin{minipage}{5in}%
                \the\cescg@abs
            \end{minipage}\par%
            \vspace*{\fill}
        \end{center}%
    \end{titlepage}%
}
%    \end{macrocode}
%  \end{macro}
%
% We cause the cover page to be generated at the end of the paper.
%    \begin{macrocode}
\AtEndDocument{\ifcescg@cover \cescg@coverpage \fi}
%    \end{macrocode}
%
% \subsection{Configurations file}
%
% Finally, we load the configuration file if it exists.
%
%    \begin{macrocode}
\InputIfFileExists{cescg.cfg}
        {\typeout{***************************************^^J%
                  * Local config file cescg.cfg used    *^^J%
                  ***************************************}}
        {}
%    \end{macrocode}
%
% \iffalse
%</class>
% \fi
%
% \iffalse ============== EXAMPLE FILE ================================
%<*sample>
% \fi
%
% \iffalse Don't include the example file in the User's Guide
% \section{Example File}
%
%    \begin{macrocode}
 %%% The document class command loads the CESCG class and sets the basic
 %%%  options.  See the user's guide for more options and their meaning.

%<*gobble>
        %% GZ: Given that the review etc. options are deprecated anyway,
        %% there is no need to have these lines here.       
\documentclass[cameraready]{cescg} % Use this for the final paper
 %deprecated - do not use this \documentclass[review]{cescg}     % Use this for a submission for review
 %deprecated - do not use this \documentclass[widereview]{cescg} % Use this for a wide spaced review paper
 %deprecated - do not use this \documentclass[preprint]{cescg}   % Use this for a preprint
 %deprecated - do not use this \preprinttext{Never appearing anywhere}  % Uncomment to change preprint info
%</gobble> 

\documentclass{cescg}[2005/11/12] % Use this for your submission and the final paper
 %%% Use e.g. this package to include figures via \includegraphics[options]{filename}.
 %%% If you don't have it in your latex installation, complain to your
 %%% sysadmin that he's not doing his job.
\usepackage{graphicx}
 %%% The following packages are somewhat system and installation dependent, and may not work
 %%% on some systems. 
 %%% If you use pdfLaTeX on Windows (part of MiKTeX), you may use package epstopdf to automatically
 %%% convert eps files to pdf. Use this line:
 %\usepackage{epstopdf}
 %%% Linux users may have to call epstopdf on the command line or configure a Makefile appropriately.

 %%% If you know how to create bounding boxes for png and jpg files, you can uncomment this
 %%% to include png files etc. directly. 
 %%% For example bmeps -b myfigure.png >myfigure.bb will do the trick in Windows.
 %%% Also, dvips needs the option "-I c3r8f" to be able to convert the bitmaps.
 %%% See the bmeps faq for more info: http://bmeps.sourceforge.net/faq.html
 %\ifx\pdfoutput\undefined \DeclareGraphicsExtensions{.eps,.png,.gif,.jpg} \fi


 %%% Now we set the fields for the title block and cover sheet
 %%% See the user's guide for more information on these items

 %% Title, author(s), and affiliation
 %%  Note that we have included both individual affiliations for each author
 %%   within the author block, as well as a common affiliation.  Only one
 %%   of these is necessary.
 %%  Footnotes to items in the title block should be created with the \thanks
 %%   command.
\title{What to Do With Money}
\author{Andreas Xnai\thanks{andreas@xnai.com} %
        \and John Jim Brown\thanks{jjb@widget.com}}
\supervisor{Molly Teacher\thanks{mt@cescg.org}}
\affiliation{Institute of Computer Graphics\\
             Whatever University of Technology\\
             Alexandria~/~Egypt}

%<*gobble>
 %% GZ: We don't need this currently.
           
 %% Paper type
\cescgcategory{survey}
\cescgformat{print}

 %% Who to contact about the paper
\contactname{Paula Vomlach}
\contactaddress{   On the Way 2130\\
                   888 88 Bratislava \\
                   Slovak Republic \\}
\contactphone{(123) 456--7890}
\contactfax{(321) 654--0987}
\contactemail{paula@pixel.area.volume.com}
%</gobble>
 %% Keywords of paper
\keywords{Bank Sector, Save Trees, Sparefroh}
% \iffalse GZ: 2005-11-28: taken out, we don't need it currently.
%<*gobble>
% estimated length of paper
\estpages{3}

 %%% GZ: Currently unused
 %%% This command suppresses printing of the cover page.  To conserve paper
 %%%  (and compile time), you should probably not comment it out until you
 %%%  are ready to send in the paper for review.
\suppresscover
%</gobble>
% \fi
 %%% Begin the paper
\begin{document}

 %%% Create the title block
\maketitle

 %%% The abstract
\begin{abstract}
  In the time being the popularity of money is permanently decreasing. 
  The question what to do with unusable money arose among the others.
  We would like to consider all aspects of this phenomenon according to the 
  present knowledge and the state of the art. We collected huge amount of 
  ideas and last ten years we sort them. Unfortunately up to now we have 
  not found two of the same category. But our effort is not at the end.
\end{abstract}

 %%% Keywords 

\keywordlist

 %%% The main text
\section{Introduction}

This is the main text of the contribution containing various information on 
the topic, which has several aspects including those mentioned in Barry Williams
book~\cite{Williams}. It is not terribly fun for the rest of us.
No one of us can say wht will be in future with money. They are becoming useless.
In this paper, we detail the results of our study.%
\footnote{Note: This is a note to the text.}

This\footnote{This is another note to the text.}
 is the main text of the contribution containing various information on 
the topic, which has several aspects including those mentioned in a book by
Karabas~\cite{Karabas}. It is not terribly fun for the rest of us.
No one of us can say what will be in future with money. They are becoming useless.
In this paper, we detail the results of our study. Compare with the literature,
e.g.~\cite{Williams,Karabas}%

In Section~\ref{sec:what}, we describe the present state of methods used in the 
 area.
In Section~\ref{sec:who}, we discuss difficulties and describe
 techniques which we subsequently developed (see Figure~\ref{fig:example}) to overcome them.
Finally, in Section~\ref{sec:why}, we examine the effects of these
 techniques on mind health of the user.

\section{Background}\label{sec:what}

Blah blah blah etc etc so on and so on and so on ad infinitum the quick brown
fox jumped over the lazy dog that came to the aid of its country.
Blah blah blah etc etc so on and so on and so on ad infinitum the quick brown
fox jumped over the lazy dog that came to the aid of its country.
Blah blah blah etc etc so on and so on and so on ad infinitum the quick brown
fox jumped over the lazy dog that came to the aid of its country.
Blah blah blah etc etc so on and so on and so on ad infinitum the quick brown
fox jumped over the lazy dog that came to the aid of its country.
Blah blah blah etc etc so on and so on and so on ad infinitum the quick brown
fox jumped over the lazy dog that came to the aid of its country.
Blah blah blah etc etc so on and so on and so on ad infinitum the quick brown
fox jumped over the lazy dog that came to the aid of its country.
 \footnote{This is one more footnote for you.}

Blah blah blah etc etc so on and so on and so on ad infinitum the quick brown
fox jumped over the lazy dog that came to the aid of its country.
Blah blah blah etc etc so on and so on and so on ad infinitum the quick brown
fox jumped over the lazy dog that came to the aid of its country.
Blah blah blah etc etc so on and so on and so on ad infinitum the quick brown
fox jumped over the lazy dog that came to the aid of its country.
Blah blah blah etc etc so on and so on and so on ad infinitum the quick brown
fox jumped over the lazy dog that came to the aid of its country.
Blah blah blah etc etc so on and so on and so on ad infinitum the quick brown
fox jumped over the lazy dog that came to the aid of its country.
Blah blah blah etc etc so on and so on and so on ad infinitum the quick brown
fox jumped over the lazy dog that came to the aid of its country.

Blah blah blah etc etc so on and so on and so on ad infinitum the quick brown
fox jumped over the lazy dog that came to the aid of its country.
Blah blah blah etc etc so on and so on and so on ad infinitum the quick brown
fox jumped over the lazy dog that came to the aid of its country.
Blah blah blah etc etc so on and so on and so on ad infinitum the quick brown
fox jumped over the lazy dog that came to the aid of its country.
Blah blah blah etc etc so on and so on and so on ad infinitum the quick brown
fox jumped over the lazy dog that came to the aid of its country.
Blah blah blah etc etc so on and so on and so on ad infinitum the quick brown
fox jumped over the lazy dog that came to the aid of its country.
Blah blah blah etc etc so on and so on and so on ad infinitum the quick brown
fox jumped over the lazy dog that came to the aid of its country.

Blah blah blah etc etc so on and so on and so on ad infinitum the quick brown
fox jumped over the lazy dog that came to the aid of its country.
Blah blah blah etc etc so on and so on and so on ad infinitum the quick brown
fox jumped over the lazy dog that came to the aid of its country.
Blah blah blah etc etc so on and so on and so on ad infinitum the quick brown
fox jumped over the lazy dog that came to the aid of its country.
Blah blah blah etc etc so on and so on and so on ad infinitum the quick brown
fox jumped over the lazy dog that came to the aid of its country.
Blah blah blah etc etc so on and so on and so on ad infinitum the quick brown
fox jumped over the lazy dog that came to the aid of its country.
Blah blah blah etc etc so on and so on and so on ad infinitum the quick brown
fox jumped over the lazy dog that came to the aid of its country.

\section{Getting Away With It}\label{sec:who}

Blah blah blah etc etc so on and so on and so on ad infinitum the quick brown
fox jumped over the lazy dog that came to the aid of its country.
Blah blah blah etc etc so on and so on and so on ad infinitum the quick brown
fox jumped over the lazy dog that came to the aid of its country.
Blah blah blah etc etc so on and so on and so on ad infinitum the quick brown
fox jumped over the lazy dog that came to the aid of its country.
Blah blah blah etc etc so on and so on and so on ad infinitum the quick brown
fox jumped over the lazy dog that came to the aid of its country.
Blah blah blah etc etc so on and so on and so on ad infinitum the quick brown
fox jumped over the lazy dog that came to the aid of its country.
Blah blah blah etc etc so on and so on and so on ad infinitum the quick brown
fox jumped over the lazy dog that came to the aid of its country.

%%% Figure inclusion. If the file extension is left away, pdfLaTeX will use
%%% png/gif/jpg/pdf versions of your figures automatically if you provide them.
%%% "Classic" LaTeX needs .eps figures with BoundingBox information.
\begin{figure}[ht]
  \centering
    \includegraphics[width=0.8\linewidth]{example}
  \caption{Example Tree}
  \label{fig:example}
\end{figure}

Blah blah blah etc etc so on and so on and so on ad infinitum the quick brown
fox jumped over the lazy dog that came to the aid of its country.
Blah blah blah etc etc so on and so on and so on ad infinitum the quick brown
fox jumped over the lazy dog that came to the aid of its country.
Blah blah blah etc etc so on and so on and so on ad infinitum the quick brown
fox jumped over the lazy dog that came to the aid of its country.
Blah blah blah etc etc so on and so on and so on ad infinitum the quick brown
fox jumped over the lazy dog that came to the aid of its country.
Blah blah blah etc etc so on and so on and so on ad infinitum the quick brown
fox jumped over the lazy dog that came to the aid of its country.%
 \footnote{You're not still reading are you? It just goes on like this.}.
Blah blah blah etc etc so on and so on and so on ad infinitum the quick brown
fox jumped over the lazy dog that came to the aid of its country.

Blah blah blah etc etc so on and so on and so on ad infinitum the quick brown
fox jumped over the lazy dog that came to the aid of its country.
Blah blah blah etc etc so on and so on and so on ad infinitum the quick brown
fox jumped over the lazy dog that came to the aid of its country.
Blah blah blah etc etc so on and so on and so on ad infinitum the quick brown
fox jumped over the lazy dog that came to the aid of its country.
Blah blah blah etc etc so on and so on and so on ad infinitum the quick brown
fox jumped over the lazy dog that came to the aid of its country.
Blah blah blah etc etc so on and so on and so on ad infinitum the quick brown
fox jumped over the lazy dog that came to the aid of its country.
Blah blah blah etc etc so on and so on and so on ad infinitum the quick brown
fox jumped over the lazy dog that came to the aid of its country.

Blah blah blah etc etc so on and so on and so on ad infinitum the quick brown
fox jumped over the lazy dog that came to the aid of its country.
Blah blah blah etc etc so on and so on and so on ad infinitum the quick brown
fox jumped over the lazy dog that came to the aid of its country.
Blah blah blah etc etc so on and so on and so on ad infinitum the quick brown
fox jumped over the lazy dog that came to the aid of its country.
Blah blah blah etc etc so on and so on and so on ad infinitum the quick brown
fox jumped over the lazy dog that came to the aid of its country.
Blah blah blah etc etc so on and so on and so on ad infinitum the quick brown
fox jumped over the lazy dog that came to the aid of its country.
Blah blah blah etc etc so on and so on and so on ad infinitum the quick brown
fox jumped over the lazy dog that came to the aid of its country.

\section{Results}\label{sec:why}

As shown in figure \ref{fig:example} and table \ref{table_1}
Blah blah blah etc etc so on and so on and so on ad infinitum the quick brown
fox jumped over the lazy dog that came to the aid of its country.
Blah blah blah etc etc so on and so on and so on ad infinitum the quick brown
fox jumped over the lazy dog that came to the aid of its country.
Blah blah blah etc etc so on and so on and so on ad infinitum the quick brown
fox jumped over the lazy dog that came to the aid of its country.
Blah blah blah etc etc so on and so on and so on ad infinitum the quick brown
fox jumped over the lazy dog that came to the aid of its country.
Blah blah blah etc etc so on and so on and so on ad infinitum the quick brown
fox jumped over the lazy dog that came to the aid of its country.
Blah blah blah etc etc so on and so on and so on ad infinitum the quick brown
fox jumped over the lazy dog that came to the aid of its country.

\begin{table}[t]
  \centering
   \begin{tabular}{|r|c|l|}
     \hline
     heading 1&heading 2&heading 3\\
     \hline
     1&2&3\\
     one&two&three\\
     \hline
   \end{tabular}
   \caption{Example Table}
  \label{table_1}
\end{table}

Blah blah blah etc etc so on and so on and so on ad infinitum the quick brown
fox jumped over the lazy dog that came to the aid of its country.
Blah blah blah etc etc so on and so on and so on ad infinitum the quick brown
fox jumped over the lazy dog that came to the aid of its country.
Blah blah blah etc etc so on and so on and so on ad infinitum the quick brown
fox jumped over the lazy dog that came to the aid of its country.
Blah blah blah etc etc so on and so on and so on ad infinitum the quick brown
fox jumped over the lazy dog that came to the aid of its country.
Blah blah blah etc etc so on and so on and so on ad infinitum the quick brown
fox jumped over the lazy dog that came to the aid of its country.
Blah blah blah etc etc so on and so on and so on ad infinitum the quick brown
fox jumped over the lazy dog that came to the aid of its country.

Blah blah blah etc etc so on and so on and so on ad infinitum the quick brown
fox jumped over the lazy dog that came to the aid of its country.
Blah blah blah etc etc so on and so on and so on ad infinitum the quick brown
fox jumped over the lazy dog that came to the aid of its country.
Blah blah blah etc etc so on and so on and so on ad infinitum the quick brown
fox jumped over the lazy dog that came to the aid of its country.
Blah blah blah etc etc so on and so on and so on ad infinitum the quick brown
fox jumped over the lazy dog that came to the aid of its country.
Blah blah blah etc etc so on and so on and so on ad infinitum the quick brown
fox jumped over the lazy dog that came to the aid of its country.
Blah blah blah etc etc so on and so on and so on ad infinitum the quick brown
fox jumped over the lazy dog that came to the aid of its country.

\section{Conclusions}

Blah blah blah etc etc so on and so on and so on ad infinitum the quick brown
fox jumped over the lazy dog that came to the aid of its country.
Blah blah blah etc etc so on and so on and so on ad infinitum the quick brown
fox jumped over the lazy dog that came to the aid of its country.
Blah blah blah etc etc so on and so on and so on ad infinitum the quick
brown fox jumped over the lazy dog that came to the aid of its country.
Blah blah blah etc etc so on and so on and so on ad infinitum the quick brown
fox jumped over the lazy dog that came to the aid of its country.
Blah blah blah etc etc so on and so on and so on ad infinitum the quick brown
fox jumped over the lazy dog that came to the aid of its country.
Blah blah blah etc etc so on and so on and so on ad infinitum the quick brown
fox jumped over the lazy dog that came to the aid of its country.

 %%% The bibliography
\bibliographystyle{plain}
\bibliography{smplbib}

\end{document}
%    \end{macrocode}
% \fi
%
% \iffalse
%</sample>
% \fi
%
%
% \iffalse
%<*bib>
% \fi
%
% \iffalse Don't include the sample bibliography in the User's Guide
% \section{Bibliography for Sample File}
%
%    \begin{macrocode}
@BOOK{Williams,
 author={Barry Williams},
 title={Money, Money, Money},
 publisher={Andrews and McMeel},
 address={Black City},
 year=1998,
}

@BOOK{Karabas,
 author={Barabas Karabas},
 title={How to Make Money out of Wood},
 publisher={Sprinter},
 year=1998,
}
%    \end{macrocode}
% \fi
%
% \iffalse
%</bib>
% \fi
%<*eps>
%<<EOF
%!PS-Adobe-2.0 EPSF-2.0
%%Title: example.eps
%%Creator: fig2dev Version 3.2 Patchlevel 3d
%%CreationDate: Sun Oct 21 16:13:08 2001
%%For: atrax@Yggdrasill (Andreas Traxler,home,,)
%%BoundingBox: 0 0 177 121
%%Magnification: 1.0000
%%EndComments
/$F2psDict 200 dict def
$F2psDict begin
$F2psDict /mtrx matrix put
/col-1 {0 setgray} bind def
/col0 {0.000 0.000 0.000 srgb} bind def
end
save
newpath 0 121 moveto 0 0 lineto 177 0 lineto 177 121 lineto closepath clip newpath
-109.5 315.1 translate
1 -1 scale

/cp {closepath} bind def
/ef {eofill} bind def
/gr {grestore} bind def
/gs {gsave} bind def
/sa {save} bind def
/rs {restore} bind def
/l {lineto} bind def
/m {moveto} bind def
/rm {rmoveto} bind def
/n {newpath} bind def
/s {stroke} bind def
/sh {show} bind def
/slc {setlinecap} bind def
/slj {setlinejoin} bind def
/slw {setlinewidth} bind def
/srgb {setrgbcolor} bind def
/rot {rotate} bind def
/sc {scale} bind def
/sd {setdash} bind def
/ff {findfont} bind def
/sf {setfont} bind def
/scf {scalefont} bind def
/sw {stringwidth} bind def
/tr {translate} bind def
/tnt {dup dup currentrgbcolor
  4 -2 roll dup 1 exch sub 3 -1 roll mul add
  4 -2 roll dup 1 exch sub 3 -1 roll mul add
  4 -2 roll dup 1 exch sub 3 -1 roll mul add srgb}
  bind def
/shd {dup dup currentrgbcolor 4 -2 roll mul 4 -2 roll mul
  4 -2 roll mul srgb} bind def
 /DrawEllipse {
        /endangle exch def
        /startangle exch def
        /yrad exch def
        /xrad exch def
        /y exch def
        /x exch def
        /savematrix mtrx currentmatrix def
        x y tr xrad yrad sc 0 0 1 startangle endangle arc
        closepath
        savematrix setmatrix
        } def

/$F2psBegin {$F2psDict begin /$F2psEnteredState save def} def
/$F2psEnd {$F2psEnteredState restore end} def

$F2psBegin
10 setmiterlimit
 0.06299 0.06299 sc
%
% Fig objects follow
%
% Polyline
7.500 slw
n 3600 4950 m 4050 4050 l 3600 3150 l
 2250 4050 l gs col0 s gr 
% Polyline
n 4500 4950 m
 4050 4050 l gs col0 s gr 
% Polyline
n 2250 3150 m
 3600 3150 l gs col0 s gr 
% Polyline
n 1800 4500 m
 2700 4950 l gs col0 s gr 
% Polyline
n 1800 4500 m
 2250 4050 l gs col0 s gr 
% Polyline
n 1800 3600 m
 2250 4050 l gs col0 s gr 
% Ellipses
n 2242 3145 47 47 0 360 DrawEllipse gs 0.00 setgray ef gr gs col-1 s gr
n 4047 4055 47 47 0 360 DrawEllipse gs 0.00 setgray ef gr gs col-1 s gr
n 3586 3148 47 47 0 360 DrawEllipse gs 0.00 setgray ef gr gs col-1 s gr
n 2250 4047 47 47 0 360 DrawEllipse gs 0.00 setgray ef gr gs col-1 s gr
n 1794 3587 47 47 0 360 DrawEllipse gs 0.00 setgray ef gr gs col-1 s gr
n 1803 4504 47 47 0 360 DrawEllipse gs 0.00 setgray ef gr gs col-1 s gr
n 2681 4928 47 47 0 360 DrawEllipse gs 0.00 setgray ef gr gs col-1 s gr
n 3580 4948 47 47 0 360 DrawEllipse gs 0.00 setgray ef gr gs col-1 s gr
n 4483 4933 47 47 0 360 DrawEllipse gs 0.00 setgray ef gr gs col-1 s gr
$F2psEnd
rs
%EOF
%</eps>
% \Finale








